\section{Introduction}
\label{sec:introduction}

\begin{table}[t]
\centering
\caption{.}
\label{tab:method_category}
\scalebox{0.63}{
\begin{tabular}{lcccc}
\toprule
\textbf{Methods} & \makecell{\textbf{Intrinsics-}\\\textbf{Free}} & \makecell{\textbf{Extrinsics-}\\\textbf{Free}} & \makecell{\textbf{View}\\\textbf{Synthesis}} & \makecell{\textbf{Direct Depth }\\\textbf{Estimation}} \\ \midrule
MVSplat~\cite{chen2024mvsplat}, pixelSplat~\cite{charatan2024pixelsplat}, ZPressor~\cite{wang2025zpressor}, Long-LRM~\cite{ziwen2025long} & \xmark & \xmark & \xmark & \xmark \\
 DepthSplat~\cite{xu2025depthsplat} & \xmark & \xmark & \xmark & \cmark\\
NoPoSplat~\cite{ye2024no}, SPFSplat~\cite{huang2025no}, EcoSplat~\cite{park2026ecosplat} & \xmark & \cmark & \cmark & \xmark\\
VGGT~\cite{wang2025vggt}, Mast3r~\cite{mast3r_eccv24}, $\pi^3$~\cite{wang2025pi3} & \cmark & \cmark  & \xmark & \cmark\\ \midrule
AnySplat~\cite{jiang2025anysplat}, WorldMirror~\cite{liu2025worldmirror}, \textbf{VFM-Splat (ours)} & \textbf{\cmark} & \textbf{\cmark} & \textbf{\cmark} & \textbf{\cmark}\\ \bottomrule
\end{tabular}}
\end{table}

% 1 about nvs and GS
The field of Novel View Synthesis (NVS) has undergone a paradigm shift with the advent of 3D Gaussian Splatting (3DGS) \cite{kerbl20233d} and its following works \cite{Wu_2024_CVPR, park2024splinegs, lin2025longsplat, bui2025mobgsmotiondeblurringdynamic, fu2024colmap, yu2024mip}, enabling high-fidelity reconstruction and real-time rendering on static dynamic or abnormal  photmetric input condition. 

% 2 about unconstrainted NVS, direct gemometry estimation (isolate), abitrary views (but can learn from sparse views)
To move beyond the limitations of time-intensive per-scene optimization, recent feed-forward NVS methods~\cite{charatan2024pixelsplat, chen2024mvsplat, xu2025depthsplat, wang2025zpressor} have been proposed to directly predict 3D scene parameters from a set of context views.
While these methods achieve instant synthesis, they predominantly rely on calibrated camera poses, which are often unavailable in "in-the-wild" scenarios. 
To handle this, pose-free feed-forward NVS methods \cite{ye2024no, huang2025no, jiang2025anysplat} jointly infer both the scene geometry and camera extrinsics for NVS.
Despite this progress, jointly estimating precise camera parameters and synthesizing novel views from sparse, uncalibrated context images remains a formidable challenge.
% why 3DVFM and why 3DVFM with small finetuning
Recently, 3D Vision Foundation Models (3DVFMs), such as and VGGT \cite{wang2025vggt}, $\pi^3$ \cite{wang2025pi3}, and DepthAnything-V3 \cite{lin2025da3}, have demonstrated remarkable zero-shot capabilities in depth and pose estimation. 
However, directly leveraging these models for NVS often yields suboptimal results, characterized by structural blurring and floating artifacts.





We identify two fundamental obstacles that hinder 3DVFMs from achieving high-fidelity NVS. 
First, 3DVFMs often lack the multi-view consistency required for dense 3D reconstruction, leading to local geometric inconsistencies and persistent "floaters" that degrade temporal stability.

% Verify the context, target training strategy. Can we just cite spfsplat?
Next, there exists a \textit{pose-geometry discrepancy} during training. 
While target poses are inferred using features from both context and target images, the scene geometry is conditioned strictly on context views to maintain generalizability. 
This asymmetric information flow leads to a misalignment between the predicted pose and the rendered 3D structure. 
When supervised by standard photometric losses, this misalignment propagates noisy gradients that force the model to average out spatial errors, resulting in blurred textures. 
% TTO is not appropriate for training loop

% Why dont we joint training like other (can ablate this, this should be broken)
To bridge this gap, we propose a novel training framework that elevates 3DVFMs from passive initialization tools to active geometric supervisors. 
Our framework introduces two key technical contributions designed to enforce global alignment and local structural integrity.
First, we present \textbf{Self-Consistent Pose Alignment (SCPA)} to align predicted target poses to the scene geometry constructed by context views and give accurate reconstruction supervision. 
By re-estimating the target poses of the rendered proxy images, we identify the geometric bias between the initial prediction and the re-estimated poses.
Then, we back-extrapolate the initial target poses using the geometric bias to be aligned with the scene geometry constructed by the context views.
This ensures that the photometric gradients are backpropagated through a pixel-aligned path, preserving the sharpness of the scene geometry.
Next, we introduce \textbf{Geometry-Guided Opacity Distillation (GGOD)} to enforce multi-view consistency. 
GGOD distills knowledge from a pre-trained two-view teacher model using an optimization strategy. 
We resolve spatial discrepancies through a geometric alignment loss while simultaneously suppressing fundamentally inconsistent floaters via a dynamic opacity ceiling. 
This ceiling responds exponentially to the relative geometric error between the student and the teacher, effectively pruning artifacts that lack multi-view support.

Experimental results on large-scale benchmarks \cite{ling2024dl3dv} demonstrate that our framework significantly outperforms other state-of-the-art NVS methods. By synergizing global photometric alignment with local geometric distillation, our approach achieves the state-of-the-art reconstruction quality, highlighting the potential of foundation models as robust active supervisors for complex 3D synthesis tasks.